% Economics Bulletin submission — "Preliminary Results" category.
%
% Formatting follows "Instructions to authors for submissions and revisions"
% (AccessEcon) to the letter. Every rule and where it is satisfied:
%
%   Seven printed pages or fewer, excluding tables, figures and references
%       (the live submission page drops "appendices" from the exclusion list
%       the 2012 PDF gave) ........................... tools/pagecount.py verifies
%   Written in English ............................... yes
%   12pt Times Roman / CM or similar ................. mathptmx, 12pt class
%   Single spaced, one-inch margins .................. geometry margin=1in
%   Sections numbered consecutively in Arabic ........ \titleformat below
%   Section headings centred, bold, 14pt ............. \titleformat below
%   Tables and figures within the text, in place ..... [H] floats, no "table
%                                                      about here" placeholders
%   No title page ................................... document opens at \section
%   No page numbers ................................. \pagestyle{empty}
%   Displayed equation numbers in parens, right ...... amsmath default
%   Footnotes sparse, superscript Arabic ............. none used
%   Tables numbered in ROMAN numerals ................ \thetable redefined
%   Short caption directly ABOVE each table .......... captionsetup position=top
%   References: surname and date in text, no comma ... "Tiebout (1956) claims"
%   References alphabetical, descending by date ...... checked
%   Reference style per the journal's examples ....... "Title" *Journal* **Vol**,
%                                                      pages.  No DOIs — the
%                                                      journal's own examples
%                                                      carry none.
%   Hyperlinks underlined but black .................. ulem \uline + hidelinks
%
\documentclass[12pt]{article}

\usepackage[margin=1in]{geometry}
\usepackage{mathptmx}                 % Times for text and maths
\usepackage[T1]{fontenc}
\usepackage[utf8]{inputenc}
\usepackage{booktabs}
\usepackage{float}
\usepackage{titlesec}
\usepackage{caption}
\usepackage{amsmath}
\usepackage[hidelinks]{hyperref}
\usepackage[normalem]{ulem}           % \uline breaks across lines

\pagestyle{empty}                     % no page numbers anywhere
\setlength{\parindent}{0pt}
\setlength{\parskip}{6pt}

% "Sections and subsections should be numbered consecutively in Arabic numerals
% (as in section 1. and subsection 1.2). Section headings should be centered and
% in bold 12pt type."  Note the journal's own example: a period after a section
% number, none after a subsection number.
% The live submission page asks for 14pt section headings (the 2012 PDF said
% 12pt). Set the size exactly rather than trusting \large, which is 14.4pt here.
\titleformat{\section}{\centering\bfseries\fontsize{14}{17}\selectfont}{\thesection.}{0.5em}{}
\titleformat{\subsection}{\centering\bfseries\normalsize}{\thesubsection}{0.5em}{}
\titlespacing*{\section}{0pt}{12pt}{6pt}
\titlespacing*{\subsection}{0pt}{10pt}{4pt}

% "Number tables consecutively with Roman numerals in order of appearance in the
% text. A short descriptive caption should be typed directly above each table."
\renewcommand{\thetable}{\Roman{table}}
\captionsetup[table]{position=top,font=small,labelfont=bf,skip=4pt}

% Table notes carry the detail that the caption must stay short enough to omit.
% \par first: without it the minipage sets on the same line as the tabular and
% its vertical centring drives a two-line note up through the bottom rule.
\newcommand{\tnote}[1]{%
  \par\addvspace{4pt}%
  \begin{minipage}{0.92\textwidth}\footnotesize\raggedright #1\end{minipage}}

\begin{document}

% No title page. Per the instructions, the first page of the submitted PDF begins
% with the first section heading; title, authors and abstract are supplied
% through the submission interface and the journal generates the title page.

\section{Introduction}

Ten states take almost all of India's foreign direct investment. Since the
Department for Promotion of Industry and Internal Trade began reporting
state-wise FDI equity figures in October 2019, those ten have absorbed close to
97 per cent of the national total (DPIIT 2024). A natural question is whether
fiscal discipline is part of what separates them from the rest, and whether it
sorts investment among them.

Two literatures speak to this and rarely cite each other. One asks how state
fiscal policy affects growth: Trivedi and Rajmal (2011) and Panda and Sahay
(2022) both report that fiscal deficits contract subsequent growth. A
debt-sustainability literature (Misra, Gupta and Trivedi 2023, and NIPFP 2024)
asks whether current borrowing paths are solvent, but stops before asking what
those paths do to the investment a state subsequently receives. NITI Aayog (2025)
ranks states on fiscal metrics in its Fiscal Health Index and validates that
ranking against no outcome at all. Internationally, Blonigen (2005) surveys the
empirical determinants of FDI and Feld, Köhler, Palhuca and Schaltegger (2024)
find that fiscal decentralisation matters for where it locates, but the
sub-national Indian case is untested.

This note reports what a ten-state, five-year panel can and cannot establish
about that question, and the answer is more useful for its negative half. States
carrying persistently lower public debt do receive more FDI, the association is
large, and five estimators drawn from three families that share no distributional
assumptions all recover it with the same sign. But the low-debt states in this
sample are the western and southern ones and the high-debt states the northern
and eastern ones; seventy-two per cent of the cross-state variance in debt lies
between those regional groups. Once regional indicators enter, the debt
coefficient is indistinguishable from zero. We report that confound as the
finding rather than as a caveat, because it bounds what any panel of this shape
can identify, and because the design that would settle the question follows
directly from it.

\section{Data}

The panel is balanced: ten states over five fiscal years, 2019--20 to 2023--24,
giving fifty state-year cells. The ten are Maharashtra, Karnataka, Gujarat, Tamil
Nadu, Haryana, Telangana, Jharkhand, Rajasthan, West Bengal and Uttar Pradesh ---
the largest FDI recipients among the states the Sixteenth Finance Commission
reports on. Delhi ranks fourth on cumulative inflow but is a union territory,
absent from the Commission's annexures, so West Bengal enters in its place.
Forty-eight of the fifty cells carry a usable FDI observation; the two that do
not belong to Uttar Pradesh, whose DPIIT series begins only in FY 2021--22 and
whose earlier cells are treated as missing rather than zero.

Four fiscal indicators, each a share of gross state domestic product, measure
discipline: gross fiscal deficit, outstanding liabilities (debt), own tax revenue
and capital expenditure. All are computed from the Sixteenth Finance Commission's
Volume II annexures (Government of India 2025), which derive from CAG-audited
State Finance Accounts. Capital expenditure is not reported directly and is
derived from the identity that the gross fiscal deficit equals the revenue
deficit plus capital outlay plus net lending, so it includes net lending; it is
therefore a noisy proxy for productive spending, and Section 4 declines to
interpret its coefficient. FDI is DPIIT's (2024) state-wise equity series and
GSDP is MoSPI's (2024) comparable 2011--12 series, from which nominal growth is
computed year on year. Literacy and urbanisation from the 2011 Census (Office of
the Registrar General and Census Commissioner 2011) enter the pooled
specification only.

Table~\ref{tab:desc} reports the range check run before any estimation. One
feature anticipates everything that follows: debt-to-GSDP spans 17.6 to 42.5 per
cent, and that spread is almost entirely cross-sectional. Decomposing each
regressor's sum of squares, 91.5 per cent of the variation in debt lies between
states and 8.5 per cent within them, with a mean pairwise rank correlation of
0.91 between any two years. GSDP growth runs the other way entirely: 98.3 per
cent of its variation is within states. Two consequences recur below. The
within-state debt coefficient is identified off very little variation and should
be expected to be fragile, and the pooled growth specification has almost no
cross-state variation to explain.

\begin{table}[H]
\centering
\caption{Descriptive statistics.}
\label{tab:desc}
\small
\begin{tabular}{@{}lrrrrr@{}}
\toprule
Variable & $N$ & Mean & Std.\ dev. & Min & Max \\
\midrule
Fiscal deficit (\% GSDP)      & 50 & 2.82  & 1.26  & $-0.7$ & 5.4 \\
Debt (\% GSDP)                & 50 & 28.82 & 6.55  & 17.6   & 42.5 \\
Own tax revenue (\% GSDP)     & 50 & 6.47  & 0.80  & 5.2    & 8.2 \\
Capital expenditure (\% GSDP) & 50 & 2.27  & 1.01  & 0.1    & 4.6 \\
FDI (Rs crore)                & 48 & 31{,}194 & 43{,}933 & 44 & 163{,}795 \\
GSDP growth (\%)              & 50 & 10.39 & 7.35  & $-4.38$ & 26.76 \\
\bottomrule
\end{tabular}
\tnote{The two missing FDI cells are Uttar Pradesh in FY 2019--20 and FY
2020--21, for which DPIIT does not report separately.}
\end{table}

\section{Method}

Two models are estimated. Model 1 takes log FDI as the dependent variable, Model
2 nominal GSDP growth, each on the four fiscal ratios and a COVID dummy:

\begin{equation}
\ln(\mathrm{FDI}_{it}) = \alpha_i + \beta_1 \mathrm{fd}_{it} + \beta_2 \mathrm{debt}_{it}
+ \beta_3 \mathrm{otr}_{it} + \beta_4 \mathrm{capex}_{it} + \beta_5 \mathrm{covid}_t + \varepsilon_{it}
\end{equation}

Both are run twice: entity fixed effects with Driscoll and Kraay (1998) standard
errors, and pooled OLS with the two Census controls. Keeping both separates the
within-state from the between-state dimension, and that distinction turns out to
carry the argument.

Ten clusters is far too few for the cluster-robust $t$-distribution to be
reliable, so every coefficient the paper relies on is re-estimated by wild
cluster bootstrap imposing the null, with Rademacher weights over the ten states
(Cameron, Gelbach and Miller 2008). That bootstrap, not the cluster-robust
$p$-value, is the arbiter throughout. We also report a Mundlak (1978) correlated
random effects specification, in which each regressor enters twice --- as the
state's five-year mean and as the annual deviation from it --- so that the
between-state and within-state coefficients are identified jointly and can be
tested against each other inside one equation.

Because the panel is short and one of its five years is an extreme common shock,
the same prediction problem is run again through LASSO (Tibshirani 1996), Elastic
Net (Zou and Hastie 2005) and SHAP values (Lundberg and Lee 2017) from a shallow
random forest (Breiman 2001), under state-blocked cross-validation. These stay
credible at forty-eight observations where instrument-heavy panel GMM estimators
do not (Roodman 2009, and Mullainathan and Spiess 2017). If a penalised linear
selector, a tree-based explainer and a panel regression all land on the same
variable, the agreement is worth more than any one $p$-value --- subject to a
qualification Section 5 makes precise.

One methodological failure is worth reporting because it is easy to repeat. An
early machine-learning specification included the Census controls and ranked
urbanisation the dominant predictor. Urbanisation is fixed within a state across
all five years, so the forest was using it to identify which state an observation
came from, re-learning the fixed effect rather than measuring agglomeration. Both
controls are excluded from the machine-learning stage for that reason.

\section{Results}

\subsection{Debt and FDI}

The fixed-effects specification ($n=48$, within $R^2 = 0.429$) and the pooled
specification ($n=48$, $R^2 = 0.643$) disagree about the sign of debt, and both
are estimated precisely enough to take seriously (Table~\ref{tab:m1}). Pooled,
debt is negative: a one-percentage-point lower debt-to-GSDP ratio is associated
with roughly 19 per cent higher FDI, since $\exp(0.176) - 1 = 0.192$. Within
states, it is positive.

\begin{table}[H]
\centering
\caption{Model 1, log FDI.}
\label{tab:m1}
\small
\begin{tabular}{@{}lrrrr@{}}
\toprule
& \multicolumn{2}{c}{FE + Driscoll-Kraay} & \multicolumn{2}{c}{Pooled OLS + controls} \\
\cmidrule(lr){2-3}\cmidrule(lr){4-5}
Variable & Coef. & $p$ & Coef. & $p$ \\
\midrule
Fiscal deficit      & $+0.566$ & 0.000 & $+0.567$ & 0.265 \\
Debt-to-GSDP        & $+0.353$ & 0.000 & $-0.176$ & 0.007 \\
Own tax revenue     & $+0.107$ & 0.684 & $+0.325$ & 0.268 \\
Capital expenditure & $-0.173$ & 0.571 & $-0.366$ & 0.008 \\
COVID year          & $-1.406$ & 0.000 & $+0.547$ & 0.161 \\
Literacy            & ---      & ---   & $+0.101$ & 0.035 \\
Urbanisation        & ---      & ---   & $+0.005$ & 0.909 \\
\bottomrule
\end{tabular}
\tnote{Literacy and urbanisation are time-invariant and absorbed by the fixed
effects, so they enter the pooled specification only.}
\end{table}

Table~\ref{tab:boot} re-estimates these under the wild cluster bootstrap, and the
picture thins considerably. The pooled debt coefficient falls to $p = 0.057$,
short of five per cent. Literacy and capital expenditure do not survive at all,
and no claim rests on either. The within-state fiscal deficit coefficient, highly
significant under Driscoll-Kraay, collapses to $p = 0.836$; no deficit-FDI claim
is made anywhere in this note.

\begin{table}[H]
\centering
\caption{Cluster-robust and wild cluster bootstrap $p$-values.}
\label{tab:boot}
\small
\begin{tabular}{@{}lrrrl@{}}
\toprule
Coefficient & Estimate & Cluster $p$ & Bootstrap $p$ & Sig.\ at 5\% \\
\midrule
M1 pooled: debt                & $-0.176$  & 0.007 & 0.057 & No \\
M1 pooled: literacy            & $+0.101$  & 0.035 & 0.235 & No \\
M1 pooled: capital expenditure & $-0.366$  & 0.008 & 0.095 & No \\
M1 within: debt                & $+0.353$  & 0.000 & 0.031 & Yes \\
M1 within: fiscal deficit      & $+0.566$  & 0.264 & 0.836 & No \\
M2 within: debt                & $+1.097$  & 0.001 & 0.020 & Yes \\
M2 within: own tax revenue     & $+4.083$  & 0.045 & 0.028 & Yes \\
M2 within: COVID dummy         & $-13.783$ & 0.000 & 0.003 & Yes \\
\bottomrule
\end{tabular}
\tnote{The bootstrap imposes the null and uses Rademacher weights over the ten
states, and is the appropriate reference distribution at this number of clusters.}
\end{table}

\subsection{Between states, not within them}

The sign reversal is not an artefact of comparing two estimators with different
assumptions. The Mundlak specification recovers it inside a single equation
(Table~\ref{tab:mundlak}), and the hypothesis that the within-state and
between-state debt coefficients are equal is rejected at $p = 0.017$ under the
same bootstrap. A between estimator fitted to the ten state means puts the
cross-state coefficient at $-0.237$, so the pooled estimate is if anything
conservative.

\begin{table}[H]
\centering
\caption{Mundlak decomposition of the debt coefficient.}
\label{tab:mundlak}
\small
\begin{tabular}{@{}lrrl@{}}
\toprule
Specification & Coef. & $p$ & What it identifies \\
\midrule
Mundlak: within-state deviation & $+0.370$ & 0.0002 & Movement around a state's own mean \\
Mundlak: state-mean term        & $-0.256$ & 0.0001 & Differences between state averages \\
Test of equality                & $0.626$  & 0.017  & The two dimensions differ \\
Between estimator               & $-0.237$ & 0.022  & Cross-state relation, direct \\
\bottomrule
\end{tabular}
\tnote{Standard errors clustered by state. The equality test is reported at its
wild cluster bootstrap $p$-value.}
\end{table}

The two coefficients answer different questions. The pooled one asks whether a
state that persistently carries less debt than its peers attracts more investment
than they do; it does. The within-state one asks whether a state's own FDI rises
when its own debt rises above its usual level. Over this window it does, and not
because borrowing attracts investors: deficits and debt rose everywhere at once
in FY 2020--21 while FDI fell everywhere, and inflows then recovered while debt
stayed at its elevated level. The within-state estimate describes the COVID
window rather than a behavioural relationship, and Section 5 shows it does not
survive a full set of year effects.

\subsection{Machine-learning agreement, and growth}

Both machine-learning families rank debt first among the five fiscal predictors
by a wide margin (Table~\ref{tab:ml}), agreeing in sign with the pooled
regression. Elastic Net is a close relative of LASSO rather than an independent
test, and we count them as one verdict.

\begin{table}[H]
\centering
\caption{Machine-learning importance, five fiscal predictors.}
\label{tab:ml}
\small
\begin{tabular}{@{}lrrr@{}}
\toprule
Variable & LASSO & Elastic Net & Mean $|$SHAP$|$ \\
\midrule
Debt-to-GSDP        & $-1.376$ & $-1.280$ & 1.029 \\
Fiscal deficit      & $+0.563$ & $+0.488$ & 0.219 \\
Capital expenditure & $-0.553$ & $-0.562$ & 0.136 \\
Own tax revenue     & $+0.039$ & $+0.098$ & 0.208 \\
COVID year          & $+0.129$ & $+0.161$ & 0.007 \\
\bottomrule
\end{tabular}
\tnote{LASSO and Elastic Net coefficients are standardised, under state-blocked
cross-validation. SHAP importance is the mean absolute SHAP value from a shallow
random forest across all forty-eight observations.}
\end{table}

The growth hypothesis fails outright. Under a single COVID dummy, within-state
debt ($+1.097$) and own tax revenue ($+4.083$) both clear five per cent under the
bootstrap. Under a full set of year fixed effects the own tax revenue coefficient
vanishes ($+0.896$, $p=0.611$) and the debt coefficient reverses sign ($-0.689$,
$p = 0.086$). Nothing in the growth equation is stable across the two time
controls, in either timing, and the pooled form has under two per cent of its
dependent variable's variation between states to explain in the first place. We
report the growth hypothesis as unsupported by this panel.

\subsection{Robustness and diagnostics}

Two objections arrive before a referee has finished the results, and both can be
answered from the panel. The first is that a five-year window containing a
pandemic may be driving everything. The second is that ten states leave any
estimate hostage to one of them. Table~\ref{tab:robust} refits the within-state
specification of Model 1 dropping, in turn, the COVID year, Jharkhand, and Uttar
Pradesh.

\begin{table}[H]
\centering
\caption{Model 1 within-state coefficients across subsamples.}
\label{tab:robust}
\small
\begin{tabular}{@{}lrrrrr@{}}
\toprule
Subsample & $n$ & Fiscal deficit & Debt & Own tax rev. & Capital exp. \\
\midrule
Baseline, all data          & 48 & $+0.566$ & $+0.353$ & $+0.107$ & $-0.173$ \\
                            &    & (0.000)  & (0.000)  & (0.684)  & (0.571)  \\
Excluding FY 2020--21       & 39 & $+0.802$ & $+0.324$ & $+0.535$ & $-0.207$ \\
                            &    & (0.017)  & (0.002)  & (0.260)  & (0.480)  \\
Excluding Jharkhand         & 43 & $-0.310$ & $+0.182$ & $+0.019$ & $+0.414$ \\
                            &    & (0.141)  & (0.015)  & (0.922)  & (0.084)  \\
Excluding Uttar Pradesh     & 45 & $+0.532$ & $+0.366$ & $+0.092$ & $-0.186$ \\
                            &    & (0.000)  & (0.000)  & (0.711)  & (0.555)  \\
\bottomrule
\end{tabular}
\tnote{Driscoll-Kraay $p$-values in parentheses.}
\end{table}

Debt is the stable coefficient: $+0.353$ at baseline, $+0.324$ without the COVID
year and $+0.366$ without Uttar Pradesh, and $+0.182$ without Jharkhand ---
smaller, but the same sign and still significant. That Uttar Pradesh changes
almost nothing is the specific check its two censored cells invite. The fiscal
deficit coefficient is not stable: dropping Jharkhand reverses its sign
altogether. This is consistent with the bootstrap verdict of
Table~\ref{tab:boot}, which already declines to credit a deficit effect, and it
is why none is claimed here.

\begin{table}[H]
\centering
\caption{Specification and diagnostic tests.}
\label{tab:diag}
\small
\begin{tabular}{@{}lll@{}}
\toprule
Test & Model 1 (log FDI) & Model 2 (growth) \\
\midrule
F-test for poolability                   & 6.741 $(p<0.0001)$ & 0.818 $(p=0.603)$ \\
Breusch-Pagan LM, pooled vs.\ RE         & 0.397 $(p=0.529)$  & 2.569 $(p=0.109)$ \\
Pesaran CD, cross-sectional dependence   & 0.019 $(p=0.985)$  & 8.397 $(p<0.0001)$ \\
Wooldridge, first-order serial corr.     & $-0.051$ $(p=0.061)$ & $-0.092$ $(p=0.021)$ \\
Modified Wald, groupwise heterosked.     & 347.33 $(p<0.0001)$ & 184.58 $(p<0.0001)$ \\
Ramsey RESET, functional form            & 1.455 $(p=0.246)$  & --- \\
\bottomrule
\end{tabular}
\tnote{The Hausman test is not computable for Model 1 and returns $p=0.177$ for
Model 2. Wooldridge tests $H_0: \rho = -0.5$.}
\end{table}

Three of these bear on how the estimates should be read. Groupwise
heteroskedasticity is present in both models, which is why every standard error
reported here is robust. Cross-sectional dependence and serial correlation are
present in the growth equation and not in the FDI equation, which is among the
reasons Section 4.3 declines to claim a growth result. And the poolability test
rejects for Model 1 but not for Model 2, which is why the FDI specification
carries state fixed effects while the growth specification gains little from
them.

\section{The regional confound}

Group the ten states by region and the debt ordering is largely a regional
ordering. Mean debt runs 20.6 per cent of GSDP in the west (Maharashtra,
Gujarat), 26.2 in the south (Karnataka, Tamil Nadu, Telangana), 33.1 in the north
(Haryana, Rajasthan, Uttar Pradesh) and 34.5 in the east (Jharkhand, West
Bengal). Seventy-two per cent of the cross-state variance in debt lies between
those four groups rather than within them. Regional indicators alone explain more
of log FDI than debt alone does, at $R^2 = 0.60$ against $0.31$.

Table~\ref{tab:region} reports what happens to each of the paper's estimates when
regional indicators are added. The pooled coefficient reverses sign and becomes
indistinguishable from zero. The Mundlak between-state term reverses. The rank
correlation collapses. LASSO promotes a western-region indicator above debt.
Within regions the ordering is not reliable either: West Bengal carries more debt
than Jharkhand and receives more investment; Rajasthan carries more than Haryana
and receives less.

\begin{table}[H]
\centering
\caption{The debt result with and without regional indicators.}
\label{tab:region}
\small
\begin{tabular}{@{}llrr@{}}
\toprule
Method & Statistic & Without region & With region \\
\midrule
Pooled OLS        & Debt coefficient       & $-0.176$ & $+0.036$ \\
Mundlak           & Between-state debt     & $-0.256$ & $+0.151$ \\
Rank correlation  & Spearman's $\rho$      & $-0.842$ & $-0.285$ \\
LASSO             & Standardised debt coef.& $-1.376$ & $-0.679$ \\
Random forest     & Mean $|$SHAP$|$, debt  & $1.029$  & $1.006$ \\
\bottomrule
\end{tabular}
\tnote{Rank correlation is Spearman's $\rho$ between five-year mean debt and
five-year total FDI across the ten states; ``with region'' removes regional means
from both series. The four regions are west, south, north and east as grouped in
the text.}
\end{table}

Only the SHAP ranking is unmoved. That is the weakest of the five as a defence,
because a regional indicator is constant within a state across all five years ---
exactly the property that made urbanisation unusable in Section 3 --- so a tree
can use it to identify states rather than to measure geography. The defence
applies to the machine-learning results and not to the pooled regression or the
rank correlation, and those are where the cross-state claim actually lives.

Two things follow. First, the convergence of five methods is not independent
evidence against this confound: they converge partly by reading the same
confounded cross-section. Second, the confound cannot be resolved by any
refinement of this sample. With ten states and four regions, regional effects
absorb three of the nine between-state degrees of freedom, and debt and region
are close to collinear by construction. A reader who prefers the regional reading
cannot be refuted from these data.

One further alternative deserves stating, because it pushes the same way. DPIIT
records equity inflows against the state in which the recipient company is
registered, not the state in which the investment is deployed. That error is not
classical: it concentrates in the states hosting corporate headquarters, which
here are Maharashtra, Karnataka and Gujarat, three of the four lowest-debt states
in the panel. A registered-office effect would therefore bias the between-state
coefficient in exactly the direction we report. The closest available test is a
market-size control, since headquarters concentration follows market size above
all: adding log GSDP to the pooled specification leaves the debt coefficient at
$-0.173$, essentially unmoved, while collapsing the literacy and capital
expenditure coefficients, which are themselves scale proxies. That does not rule
the channel out --- a control for size is not a control for where a firm chooses
to register --- but it does mean the concern cannot be carried by state size
alone.

That is a statement about the design rather than about the question, and the
distinction matters because the Commission's annexures already contain the fiscal
side of the remedy. The Sixteenth Finance Commission reports all four indicators
for twenty-eight states and DPIIT reports a cumulative inflow for each. Across
the ten states studied here, 72 per cent of the cross-sectional variance in debt
lies between regions and the mean within-region spread is 6.7 percentage points.
Across twenty-six states with both series available, the between-region share
falls to 35 per cent and the mean within-region spread more than doubles to 15.1
points, with six regions averaging over four states each rather than four regions
averaging two and a half. A cross-section on those twenty-six states points the
same way as this panel: debt alone gives $-0.299$ ($p = 0.004$), and with
regional indicators added it holds at $-0.143$ ($p = 0.056$) instead of collapsing
to zero as it does here. We do not report that as a result --- the wider set has
no GSDP series in our hands and so no market-size control, and a cumulative total
over a different window is not the annual outcome modelled here. We report it as
evidence that the identifying design exists and is within reach.

\section{Conclusion}

Foreign direct investment in India is more likely where a state holds its debt
down relative to its own economy, but only between states, not within a state
from one year to the next. Five estimators from three families agree on the sign,
and the between-within split survives inside a single equation. What none of them
can do is separate that fiscal signal from a regional one. Ten states across four
regions do not contain the variation required, and the honest report of this
panel is the identification problem rather than the coefficient.

Two implications follow for practice. For a state finance department, the reading
consistent with these estimates is that fiscal position operates as a structural,
between-state signal rather than a lever: a state that has held debt below 25 per
cent of GSDP for five consecutive years is in a different competitive tier from
one that has not, and no single budget year moves that. For an index builder ---
NITI Aayog (2025) being the immediate case --- the result is a caution. The
composite ranking constructed here is outperformed by debt alone on every outcome
tested, and a composite validated against no outcome cannot know whether it is
measuring anything.

Three limits bound all of this beyond the identification problem itself. The
panel is short and one of its five years is an extreme common shock, so the
within-state estimates describe that window as much as any behavioural
relationship. Debt is the on-budget figure from the audited accounts and excludes
off-budget guarantees, which several of these states use. And association is all
a five-year panel of ten states can support: reverse causality, in which
anticipated investment relaxes a state's borrowing constraint rather than the
reverse, is not tested here and cannot be ruled out.

The extension that would settle the question is specific rather than aspirational:
a panel wide enough to hold region fixed while debt still varies within it.
Twenty-six states clear that bar where ten do not.

\section*{Data and code availability}

The dataset, the analysis code, and the full-length working paper from which this
note is drawn are openly available at

\begin{center}
\href{https://github.com/fbivinay/Fiscal-Discipline-and-Investment-Outcomes}%
{\uline{https://github.com/fbivinay/Fiscal-Discipline-and-Investment-Outcomes}}
\end{center}

The repository holds the audited source series, the fifty-row analysis panel
behind every estimate reported here, and every analysis script.

\section*{References}
\small

\hangindent=0.5in\hangafter=1
Blonigen, B.A. (2005) ``A review of the empirical literature on FDI determinants'' \textit{Atlantic Economic Journal} \textbf{33}, 383-403.

\hangindent=0.5in\hangafter=1
Breiman, L. (2001) ``Random forests'' \textit{Machine Learning} \textbf{45}, 5-32.

\hangindent=0.5in\hangafter=1
Cameron, A.C., J.B. Gelbach and D.L. Miller (2008) ``Bootstrap-based improvements for inference with clustered errors'' \textit{Review of Economics and Statistics} \textbf{90}, 414-427.

\hangindent=0.5in\hangafter=1
Department for Promotion of Industry and Internal Trade (DPIIT) (2024) \textit{State-wise FDI Equity Inflows from October 2019}, Ministry of Commerce and Industry: New Delhi.

\hangindent=0.5in\hangafter=1
Driscoll, J.C. and A.C. Kraay (1998) ``Consistent covariance matrix estimation with spatially dependent panel data'' \textit{Review of Economics and Statistics} \textbf{80}, 549-560.

\hangindent=0.5in\hangafter=1
Feld, L.P., E.A. Köhler, L. Palhuca and C.A. Schaltegger (2024) ``Fiscal federalism and foreign direct investment: an empirical analysis'' \textit{The World Economy} \textbf{47}, 2287-2331.

\hangindent=0.5in\hangafter=1
Government of India (2025) \textit{Report of the Sixteenth Finance Commission, Volume II: Annexures}, Ministry of Finance: New Delhi.

\hangindent=0.5in\hangafter=1
Lundberg, S.M. and S.-I. Lee (2017) ``A unified approach to interpreting model predictions'' \textit{Advances in Neural Information Processing Systems} \textbf{30}, 4765-4774.

\hangindent=0.5in\hangafter=1
Ministry of Statistics and Programme Implementation (MoSPI) (2024) \textit{State-wise Gross State Domestic Product, 2011-12 Series}, Government of India: New Delhi.

\hangindent=0.5in\hangafter=1
Misra, S., K. Gupta and P. Trivedi (2023) ``Sub-national government debt sustainability in India: an empirical analysis'' \textit{Macroeconomics and Finance in Emerging Market Economies} \textbf{16}, 57-79.

\hangindent=0.5in\hangafter=1
Mullainathan, S. and J. Spiess (2017) ``Machine learning: an applied econometric approach'' \textit{Journal of Economic Perspectives} \textbf{31}, 87-106.

\hangindent=0.5in\hangafter=1
Mundlak, Y. (1978) ``On the pooling of time series and cross section data'' \textit{Econometrica} \textbf{46}, 69-85.

\hangindent=0.5in\hangafter=1
National Institute of Public Finance and Policy (NIPFP) (2024) ``How much debt is optimal for the major Indian states? Economic growth vs. debt sustainability'' NIPFP working paper number 411.

\hangindent=0.5in\hangafter=1
NITI Aayog (2025) \textit{Fiscal Health Index 2025}, Government of India: New Delhi.

\hangindent=0.5in\hangafter=1
Office of the Registrar General and Census Commissioner (2011) \textit{Census of India 2011, Primary Census Abstract}, Ministry of Home Affairs: New Delhi.

\hangindent=0.5in\hangafter=1
Panda, M. and S. Sahay (2022) ``Determinants of economic growth across states in India'' in \textit{Perspectives on Inclusive Policies for Development in India} by S.R. Hashim, R. Mukherji and B. Mishra, Eds., Springer: Singapore, 133-152.

\hangindent=0.5in\hangafter=1
Roodman, D. (2009) ``How to do xtabond2: an introduction to difference and system GMM in Stata'' \textit{Stata Journal} \textbf{9}, 86-136.

\hangindent=0.5in\hangafter=1
Tibshirani, R. (1996) ``Regression shrinkage and selection via the lasso'' \textit{Journal of the Royal Statistical Society, Series B} \textbf{58}, 267-288.

\hangindent=0.5in\hangafter=1
Trivedi, P. and Rajmal (2011) ``Growth effects of fiscal policy of Indian states'' \textit{Millennial Asia} \textbf{2}, 141-162.

\hangindent=0.5in\hangafter=1
Zou, H. and T. Hastie (2005) ``Regularization and variable selection via the elastic net'' \textit{Journal of the Royal Statistical Society, Series B} \textbf{67}, 301-320.

\end{document}
